\chapter*{Résumé}
\thispagestyle{empty}

Dans un contexte industriel où la sécurité et la digitalisation constituent des enjeux majeurs, l'Office Chérifien des Phosphates (OCP) s'engage dans une démarche d'amélioration continue de ses processus de gestion des accès et de sûreté industrielle. Ce projet de stage s'inscrit dans cette dynamique en contribuant au développement et à l'optimisation du système JLHUBSECURITY chez OCP Maintenance Solutions.

\vspace{0.5cm}

Les processus traditionnels de gestion des accès au site de Jorf Lasfar présentaient plusieurs limitations : saisies manuelles répétitives sources d'erreurs, temps d'attente prolongés aux heures de pointe, absence de traçabilité numérique et risque de perte des pièces d'identité des visiteurs. Ces dysfonctionnements impactaient l'efficacité opérationnelle et la satisfaction des utilisateurs.

\vspace{0.5cm}

 L'objectif principal consistait à digitaliser et moderniser le système de gestion des accès en développant une plateforme web intégrée permettant de dématérialiser les demandes, automatiser les workflows de validation, générer des badges QR code sécurisés et améliorer l'expérience utilisateur.
 
\vspace{0.5cm}

Le projet a été mené selon une approche agile avec analyse des besoins utilisateurs, étude des processus existants, conception d'une architecture technique robuste et développement d'une solution modulaire. La stack technique retenue comprend Laravel pour le backend et Vue.js pour le frontend, garantissant performance et maintenabilité.

\vspace{0.5cm}

La solution développée intègre plusieurs modules fonctionnels : gestion des visiteurs avec génération automatique de QR codes, système de validation multi-niveaux, tableau de bord de pilotage, gestion des véhicules et du matériel, ainsi qu'un module spécialisé pour la ferraille. L'interface responsive assure une utilisation optimale sur tous les supports.

\vspace{0.5cm}

 L'implémentation du système JLHUBSECURITY a permis de réduire significativement les temps de traitement des demandes, d'éliminer les erreurs de saisie manuelle, d'améliorer la traçabilité des accès et d'offrir une meilleure expérience utilisateur. Les fonctionnalités de reporting et d'audit renforcent le contrôle et la supervision des activités.
 
\vspace{0.5cm}

Ce projet démontre l'importance de la transformation digitale dans l'industrie et contribue à positionner OCP comme un leader en matière d'innovation technologique appliquée à la sûreté industrielle.


